\documentclass[10pt,a4paper]{article}
\usepackage[margin=20mm,headheight=15pt,headsep=7mm,footskip=10mm]{geometry}
\usepackage{fontspec,amsmath,amssymb,mathtools,microtype}
\setmainfont{lmroman10-regular.otf}[BoldFont=lmroman10-bold.otf,ItalicFont=lmroman10-italic.otf,BoldItalicFont=lmroman10-bolditalic.otf]
\setsansfont{lmsans10-regular.otf}[BoldFont=lmsans10-bold.otf,ItalicFont=lmsans10-oblique.otf,BoldItalicFont=lmsans10-boldoblique.otf]
\setmonofont[Scale=.86]{lmmono10-regular.otf}[ItalicFont=lmmono10-italic.otf]
\usepackage{ragged2e}\usepackage{xcolor,booktabs,longtable,array,graphicx,enumitem,fancyhdr,needspace,titlesec,hyperref,xurl}
\definecolor{AtlasNavy}{HTML}{18334C}\definecolor{AtlasTeal}{HTML}{116D73}\definecolor{AtlasGray}{HTML}{52616D}
\hypersetup{colorlinks=true,linkcolor=AtlasTeal,urlcolor=AtlasTeal,pdfcreator={Complexity Atlas},unicode=true}
\urlstyle{same}\setlength{\parindent}{0pt}\setlength{\parskip}{4.8pt}\setlength{\emergencystretch}{3em}
\renewcommand{\baselinestretch}{1.08}\setcounter{secnumdepth}{0}\setlist{nosep,leftmargin=1.5em}
\titleformat{\section}{\Large\bfseries\sffamily\color{AtlasNavy}\raggedright}{}{0em}{}
\titleformat{\subsection}{\large\bfseries\sffamily\color{AtlasTeal}\raggedright}{}{0em}{}
\titleformat{\subsubsection}{\normalsize\bfseries\color{AtlasNavy}\raggedright}{}{0em}{}
\titlespacing*{\section}{0pt}{14pt}{6pt}\titlespacing*{\subsection}{0pt}{10pt}{5pt}
\pagestyle{fancy}\fancyhf{}\fancyhead[L]{\small\sffamily\color{AtlasGray}COMPLEXITY ATLAS}
\fancyhead[R]{\small\sffamily\color{AtlasGray}ENGLISH EDITION}\fancyfoot[C]{\small\thepage}\renewcommand{\headrulewidth}{.3pt}
\newcommand{\AtlasMath}[1]{\par\begingroup\centering\setbox0=\hbox{$\displaystyle #1$}\ifdim\wd0>1.4\linewidth\typeout{ATLAS WIDE FORMULA: \the\wd0}\fi\ifdim\wd0>\linewidth\resizebox{\linewidth}{!}{\box0}\else\box0\fi\par\endgroup}
\newcommand{\AtlasInlineMath}[1]{\begingroup\setbox0=\hbox{$#1$}\ifdim\wd0>.4\linewidth\endgroup\AtlasMath{#1}\else\endgroup$#1$\fi}
\pdfstringdefDisableCommands{\def\AtlasInlineMath#1{#1}}
\DeclareMathOperator{\polylog}{polylog}\DeclareMathOperator{\dist}{dist}\DeclareMathOperator{\diam}{diam}
\DeclareMathOperator{\rank}{rank}\DeclareMathOperator{\tr}{tr}\DeclareMathOperator{\diag}{diag}
\DeclareMathOperator{\prox}{prox}\DeclareMathOperator{\KL}{KL}\DeclareMathOperator{\Var}{Var}
\providecommand{\eps}{\varepsilon}\providecommand{\R}{\mathbb R}\providecommand{\E}{\mathbb E}
\providecommand{\mathbbm}[1]{\mathbf{#1}}\providecommand{\mathds}[1]{\mathbf{#1}}
\providecommand{\llbracket}{[\![}\providecommand{\rrbracket}{]\!]}
\begin{document}
\raggedright

{\LARGE\bfseries\sffamily\color{AtlasNavy} Stochastic bilevel gradients with two-point coupling within one fresh sample\par}\vspace{5pt}
\textbf{Setting ID.} \texttt{BL\_\allowbreak{}SPEC\_\allowbreak{}two\_\allowbreak{}point\_\allowbreak{}coupled}

Bilevel optimization

Solve the bilevel problem \AtlasInlineMath{\min_x\Phi(x)}, where \AtlasInlineMath{\Phi(x)=f(x,y^*(x))} and \AtlasInlineMath{y^*(x)=\operatorname*{argmin}_y g(x,y)} is the unique global lower-level solution. The upper- and lower-level variables are \AtlasInlineMath{x\in\mathbb{R}^{d_x}} and \AtlasInlineMath{y\in\mathbb{R}^{d_y}}. Seek approximate stationarity of the reduced objective under the declared criterion and guarantee; assumptions on \AtlasInlineMath{f}, \AtlasInlineMath{g}, and \AtlasInlineMath{\Phi} are specified separately.

\textbf{Status.} Bounds missing

\textbf{TASK tags.} task family: Bilevel

\textbf{MODEL tags.} convexity: Lower objective strongly convex in its lower variable; initialization: Bounded initial inner distance; initialization: Initial gap

\textbf{ORACLE tags.} interface family: Two point fresh

\textbf{ALGORITHMS tags.} restriction: No additional span restriction explicitly imposed

\textbf{OUTPUT tags.} postprocessing: Unqueried output explicitly allowed

\textbf{METRIC tags.} criterion: Reduced-objective gradient norm

\textbf{GUARANTEE tags.} aggregation: Per instance; mode: Expectation

\textbf{RESOURCE tags.} accounting: Joint coordinates retained; budget: Fixed worst-case budget; measure: N pair f; measure: N pair g

\needspace{7\baselineskip}\section{1. Setting specification}
\needspace{4\baselineskip}\subsection{Task}
Solve the bilevel problem \AtlasInlineMath{\min_x\Phi(x)}, where \AtlasInlineMath{\Phi(x)=f(x,y^*(x))} and \AtlasInlineMath{y^*(x)=\operatorname*{argmin}_y g(x,y)} is the unique global lower-level solution. The upper- and lower-level variables are \AtlasInlineMath{x\in\mathbb{R}^{d_x}} and \AtlasInlineMath{y\in\mathbb{R}^{d_y}}. Seek approximate stationarity of the reduced objective under the declared criterion and guarantee; assumptions on \AtlasInlineMath{f}, \AtlasInlineMath{g}, and \AtlasInlineMath{\Phi} are specified separately.

\needspace{4\baselineskip}\subsection{Admissible model}
The function \AtlasInlineMath{g(x,\cdot)} is \AtlasInlineMath{\mu_y}-strongly convex. The functions \AtlasInlineMath{\nabla f,\nabla g} are jointly \AtlasInlineMath{L_f,L_g}-Lipschitz; \AtlasInlineMath{\|\nabla_yf\|\le C_f}, and \AtlasInlineMath{g_{xy},g_{yy}} is jointly \AtlasInlineMath{\rho}-Lipschitz. Also \AtlasInlineMath{\Phi(x_0)-\inf\Phi\le\Delta} and \AtlasInlineMath{\|y_0-y^*(x_0)\|\le R_y}.

Population gradients are unbiased, with variances bounded by \AtlasInlineMath{\sigma_f^2,\sigma_g^2}. If \AtlasInlineMath{L_{\mathrm{ms}}<\infty}, additionally require the common-sample condition \AtlasInlineMath{\mathbb{E}\|G_g(x,y_1;\xi)-G_g(x,y_2;\xi)\|^2\le L_{\mathrm{ms}}^2\|y_1-y_2\|^2}; \AtlasInlineMath{L_{\mathrm{ms}}=\infty} means this condition is absent. It applies only in the lower-variable direction at fixed \AtlasInlineMath{x}.

\needspace{4\baselineskip}\subsection{Oracle interface}
The algorithm specifies at most two points \AtlasInlineMath{z_1,z_2} before sampling. Independent upper- and lower-level seeds \AtlasInlineMath{\zeta} and \AtlasInlineMath{\xi} are then drawn, and the oracle returns \AtlasInlineMath{(G_f(z_j;\zeta),G_g(z_j;\xi))} at each requested point. The two points share the current batch seeds; different batches use fresh seeds. No persistent sample handle is provided.

Queries may use the full Euclidean space. Budgets, dimensions, and initial points are public. No additional localization information about \AtlasInlineMath{y^*(x)} is supplied.

\needspace{4\baselineskip}\subsection{Output}
Return any measurable \AtlasInlineMath{\widehat{x}\in\mathbb{R}^{d_x}}, possibly an unqueried point.

\needspace{4\baselineskip}\subsection{Loss or accuracy criterion}
\AtlasInlineMath{\|\nabla\Phi(\widehat{x})\|^2}.

\needspace{4\baselineskip}\subsection{Resource accounting}
Each function-point-sample gradient evaluation costs one unit. Querying both levels at two points costs four units. The fixed worst-case budget includes all inner, outer, and initialization calls.

\needspace{4\baselineskip}\subsection{Algorithm views and success guarantees}
Arbitrary randomized adaptive algorithms are allowed, but both points in a batch must be fixed before its seed is drawn. Queries, stopping, and postprocessing may be arbitrary measurable functions; outputs may be unqueried, randomly selected historical points, or admissible averages. No span or zero-respecting restriction is imposed. Queries are sequential, and each response in a batch is charged. Internal arithmetic, linear algebra, and storage are free; internal work, memory, and communication are not budgeted.

For every instance, \AtlasInlineMath{\mathbb E\|\nabla\Phi(\widehat x)\|^2\le\varepsilon^2}, with expectation over all randomness.

\needspace{7\baselineskip}\section{2. Notation and parameters}
\AtlasInlineMath{x,y}: The upper- and lower-level variables. \textbf{Domain:} \AtlasInlineMath{x\in\mathbb R^{d_x},y\in\mathbb R^{d_y}}.

\AtlasInlineMath{f,g}: The original upper- and lower-level real-valued functions. \textbf{Domain:} \AtlasInlineMath{f,g:\mathbb R^{d_x}\times\mathbb R^{d_y}\to\mathbb R}.

\AtlasInlineMath{\widehat x,x_0,y_0}: Output outer point, public outer initial point, and public lower-level initial point \textbf{Domain:} Each lies in its corresponding variable space.

\AtlasInlineMath{\Phi}: The reduced objective. \textbf{Domain:} Defined in the task specification.

\AtlasInlineMath{y^*}: Mapping to the unique lower-level global minimizer \textbf{Domain:} Constraints specified under TASK.

\AtlasInlineMath{\nabla,\|\cdot\|}: Euclidean gradient and vector norm; matrix and higher-derivative norms are the induced Euclidean operator norms

\AtlasInlineMath{\mathbb E,\Pr,\inf,\sup}: Expectation, probability, infimum, and supremum, with aggregation order and randomness specified by the guarantee.

\AtlasInlineMath{\mathcal O,\Omega,\Theta,\widetilde{\mathcal O}}: Asymptotic upper, lower, matching-order, and declared-logarithm upper bounds; fixed parameters are specified in each evidence record

\AtlasInlineMath{c,C}: Positive absolute constants or functions of fixed budgets, with dependencies specified by the corresponding result \textbf{Domain:} \AtlasInlineMath{0<c,C<\infty}.

\AtlasInlineMath{\varepsilon}: Target accuracy. Whether error is squared is determined by the metric and guarantee formulas. \textbf{Domain:} \AtlasInlineMath{\varepsilon>0}.

\AtlasInlineMath{d_x}: Upper-level dimension \textbf{Domain:} Positive integer.

\AtlasInlineMath{d_y}: Lower-level dimension \textbf{Domain:} Positive integer.

\AtlasInlineMath{L_f}: Joint-variable Lipschitz upper bound for the complete population upper-level gradient \textbf{Domain:} \AtlasInlineMath{L_f\ge0}.

\AtlasInlineMath{L_g}: A joint-variable Lipschitz bound on the population lower-level full gradient. \textbf{Domain:} \AtlasInlineMath{L_g\ge\mu_y>0}.

\AtlasInlineMath{C_f}: A uniform norm bound on the population upper function’s lower-variable gradient block. \textbf{Domain:} \AtlasInlineMath{C_f\ge0}.

\AtlasInlineMath{\|\nabla_yf(x,y)\|\le C_f}

\AtlasInlineMath{\mu_y}: The uniform strong-convexity modulus in the lower-level variable. \textbf{Domain:} \AtlasInlineMath{\mu_y>0}.

\AtlasInlineMath{\rho}: Common joint Lipschitz upper bound for the lower function's cross-Hessian block and lower-variable Hessian block \textbf{Domain:} \AtlasInlineMath{\rho\ge0}.

\AtlasInlineMath{\Delta}: Public upper bound on the initial reduced-objective gap \textbf{Domain:} \AtlasInlineMath{\Delta\ge0}.

\AtlasInlineMath{\Phi(x_0)-\inf_x\Phi(x)\le\Delta}

\AtlasInlineMath{R_y}: Public initial lower-level distance bound \textbf{Domain:} \AtlasInlineMath{R_y\ge0}.

\AtlasInlineMath{\|y_0-y^*(x_0)\|\le R_y}

\AtlasInlineMath{L_{\rm ms}}: A mean-square Lipschitz bound on gradient differences at two lower-level points using the same lower-level sample. \textbf{Domain:} The parameter satisfies \AtlasInlineMath{0\le L_{\mathrm{ms}}\le\infty}; infinity means that this condition is not imposed.

\AtlasInlineMath{\sigma_f}: Conditional standard-deviation bound for the full stochastic upper-level gradient \textbf{Domain:} \AtlasInlineMath{\sigma_f\ge0}.

\AtlasInlineMath{\sigma_g}: Conditional standard-deviation bound for the full stochastic lower-level gradient \textbf{Domain:} \AtlasInlineMath{\sigma_g\ge0}.

\AtlasInlineMath{g_{xy},g_{yy}}: Mixed lower-level Hessian block and lower-variable Hessian block \textbf{Domain:} \AtlasInlineMath{g_{xy}\in\mathbb R^{d_x\times d_y},g_{yy}\in\mathbb R^{d_y\times d_y}}.

\AtlasInlineMath{g_{xy}=\mathrm D_y\nabla_xg,\quad g_{yy}=\mathrm D_y\nabla_yg}

\AtlasInlineMath{G_f,G_g,\mathcal H}: Upper- and lower-level stochastic full-gradient responses and the pre-draw history. \textbf{Domain:} \AtlasInlineMath{G_f,G_g\in\mathbb R^{d_x+d_y}}.

\AtlasInlineMath{G_{f,x},G_{f,y},G_{g,x},G_{g,y}}: Outer and lower blocks of the corresponding full stochastic gradient \textbf{Domain:} Outer-block dimension d\_x and lower-block dimension d\_y.

\AtlasInlineMath{z_1,z_2}: At most two joint query points specified simultaneously before one sample draw. \textbf{Domain:} \AtlasInlineMath{z_j=(x_j,y_j)}.

\AtlasInlineMath{\xi,\zeta}: Lower- and upper-level seeds, in the source order used here. They are independent.

\AtlasInlineMath{n,T,M}: The source algorithm’s predetermined batch and inner/outer iteration budgets. \textbf{Domain:} Role fixed by the evidence's call-accounting convention.

\AtlasInlineMath{N,Q}: This is the total call budget under the stated resource projection. If only one symbol is used, the two names are aliases for that same budget. \textbf{Domain:} RESOURCE specifies which operations are summed; no additional unit cost is introduced.

\needspace{12\baselineskip}\section{3. Complexity bracket}
\needspace{4\baselineskip}\subsection{Compared quantities}
\needspace{4\baselineskip}\subsubsection{Q1 - source scoped metric resource statement}
\textbf{Metric contract.} E||grad Phi(x\_hat)||\textasciicircum{}2<=epsilon\textasciicircum{}2.

\textbf{Resource contract.} vector: N\_pair\_f; N\_pair\_g; projection: Q counts each actual point/sample gradient separately; two points and two functions cost four. Fixed batch size two makes this a constant-factor mapping of the source N-simultaneous interface.; budget: Fixed outer/inner horizons and batches; all initialization and inner calls included.

This is the native source metric/resource guarantee. Joint precision, block size, runtime and oracle components stay separate; do not infer a scalar optimal query complexity.

\needspace{4\baselineskip}\subsubsection{Q2 - source scoped metric resource statement}
\textbf{Metric contract.} E||grad Phi(x\_hat)||\textasciicircum{}2<=epsilon\textasciicircum{}2.

\textbf{Resource contract.} vector: N\_pair\_f; N\_pair\_g; projection: Q counts each actual point/sample gradient separately; two points and two functions cost four. Fixed batch size two makes this a constant-factor mapping of the source N-simultaneous interface.; budget: Fixed outer/inner horizons and batches; all initialization and inner calls included.

This is the native source metric/resource guarantee. Joint precision, block size, runtime and oracle components stay separate; do not infer a scalar optimal query complexity.

\needspace{4\baselineskip}\subsection{R1 - BM\_BOUND\_coupled\_variance\_upper source window}
\textbf{Region.} No finite mean-square smoothness budget is required; \AtlasInlineMath{L_{\mathrm{ms}}=\infty} is allowed. Regularity, variance, initial-gap, and inner-distance budgets are fixed, and \AtlasInlineMath{\varepsilon} is sufficiently small.

\textbf{Quantity.} Q2

\textbf{Resource.} Joint resource budget

\textbf{Status.} Lower bound missing

\textbf{Lower bound.} No admitted bound for this comparison

\textbf{Upper bound.} B1

\needspace{4\baselineskip}\subsection{R2 - BM\_BOUND\_coupled\_mss\_upper source window}
\textbf{Region.} \AtlasInlineMath{L_{\mathrm{ms}}<\infty}. Fix regularity, variance, initial-gap, and inner-distance budgets, with \AtlasInlineMath{\varepsilon} sufficiently small.

\textbf{Quantity.} Q1

\textbf{Resource.} Joint resource budget

\textbf{Status.} Lower bound missing

\textbf{Lower bound.} No admitted bound for this comparison

\textbf{Upper bound.} B2

\needspace{7\baselineskip}\section{4. Bounds and supporting evidence}
The following conditions apply to every bound in this profile, in addition to each bound's stated region and algorithm scope.

\textbf{Oracle mapping.} In the source branch \AtlasInlineMath{r=\infty}, each fresh coupled draw is evaluated at no more than two points.

\textbf{Resource mapping.} Each native joint batch is implemented by a constant number of separately charged point-and-function gradient responses.

\needspace{4\baselineskip}\subsection{B1 - Penalty method with coupled lower iterates}
\textbf{Admission.} Included in the bracket

\textbf{Direction.} upper

\textbf{Witness or lower-bound scope.} Randomized witness

\AtlasMath{\mathcal{O}(\varepsilon^{-6})}
\textbf{Sources.} \href{https://raw.githubusercontent.com/mlresearch/v235/main/assets/kwon24b/kwon24b.pdf}{[1]}

\textbf{Constants.} Hold the stated regularity and initialization budgets fixed. Joint optimality is not claimed for parameters whose dependence is untracked.

\textbf{Source location.} \href{https://raw.githubusercontent.com/mlresearch/v235/main/assets/kwon24b/kwon24b.pdf}{[1]} pdf pages: 1; 2; 5; 6; 15; 16; 21; 22; theorem or section: Equations (2)-(4); Definition 1.1; Oracle Classes; Theorems 3.1-3.2; Eq.(86); theorem: Equations (2)-(4); Definition 1.1; Oracle Classes; Theorems 3.1-3.2; Eq.(86)

\needspace{4\baselineskip}\subsection{B2 - Penalty method with coupled lower iterates}
\textbf{Admission.} Included in the bracket

\textbf{Direction.} upper

\textbf{Witness or lower-bound scope.} Randomized witness

\AtlasMath{\mathcal{O}(\varepsilon^{-4})}
\textbf{Sources.} \href{https://raw.githubusercontent.com/mlresearch/v235/main/assets/kwon24b/kwon24b.pdf}{[1]}

\textbf{Constants.} Hold the stated regularity and initialization budgets fixed. Joint optimality is not claimed for parameters whose dependence is untracked.

\textbf{Source location.} \href{https://raw.githubusercontent.com/mlresearch/v235/main/assets/kwon24b/kwon24b.pdf}{[1]} pdf pages: 1; 2; 5; 6; 15; 16; 21; 22; theorem or section: Equations (2)-(4); Definition 1.1; Oracle Classes; Theorems 3.1-3.2; Eq.(86); theorem: Equations (2)-(4); Definition 1.1; Oracle Classes; Theorems 3.1-3.2; Eq.(86)

\needspace{7\baselineskip}\section{References}
[1] Jeongyeol Kwon; Dohyun Kwon; Hanbaek Lyu. On The Complexity of First-Order Methods in Stochastic Bilevel Optimization. 2024. PMLR235 ICML2024 publisher PDF.

\url{https://raw.githubusercontent.com/mlresearch/v235/main/assets/kwon24b/kwon24b.pdf}

\end{document}
